% !TEX root = ../HDR.tex
%%%%%%%%%%%%%%%%%%%%%%%%%%%%%%%%%%%%%%%%%%%%%%%%%%%%%%%%%%%%%%%%%%%%%%%%%%%%%%%%%%
\chapter{Discussion}



\section{Critique de la discretisation temporelle}

Les mécanismes de schéma sont basées sur des cycles d'interaction qui ne permettre qu'une prise en compte d'une temporalité discrète. 
On peut ce demander si cette limitation n'est pas trop contraignante pour permettre de générer des comportements naturels. 

Un contre argument est qu'on peut se fixer des schémas primitifs arbitrairement cours. 
Il faut cependant garder à l'esprit que le mécanisme de schéma permet de monter dans des échelles de temps de plus en plus longues. 




\section{Critique du réductionnisme}

Est ce que la théorie de Piaget est trop réductionniste et parvient elle à capture les fondements de la conscience ? 

C'est, à mon sens, la critique principale de cette approche et je n'ai pas la réponse. 